\section{Robustness}\label{sec:robustness}
We now apply the tube bounds from the previous section to the classification setting introduced in Section~\ref{sec:pfaffian}. Recall that a classifier $F\colon \R^n \to \R^m$ induces decision regions $C_j$ separated by the decision boundary $\Sigma = \bigcup_{i<j} V_{ij}$, where $V_{ij} = \mathcal{Z}(g_{ij})$ and $g_{ij} = F_j - F_i$. For $x\in \R^n\setminus \Sigma$, the distance to misclassification is
\begin{equation*}
  \Delta(x) := \dist(x, \Sigma) = \inf\{\|x - y\| \colon y \in \Sigma\}.
\end{equation*}

It is natural to measure the distance to misclassification relative to the size of the input. In analogy to the theory of conditioning in numerical analysis and optimisation~\cite{burgisser2013condition}, we define the condition of a classifier.

\begin{definition}\label{def:condition}
The \emph{condition number} of the classifier $F$ at $x\in \R^n\setminus \Sigma$ is
\begin{equation*}
  \mC(x) := \frac{\|x\|}{\Delta(x)}.
\end{equation*}
More generally, the \emph{local condition number} centred at $p\in \R^n$ is
\begin{equation*}
  \mC_p(x) := \frac{\|x - p\|}{\Delta(x)}.
\end{equation*}
\end{definition}

The condition number $\mC(x)$ is the reciprocal of the relative distance to the decision boundary, so that $\mC(x) > t$ means the classification of $x$ can be changed by a perturbation of relative size less than $1/t$. This is entirely analogous to the role of the condition number in numerical analysis, where the ``ill-posed problems'' correspond to the decision boundary~$\Sigma$.

\begin{remark}[Weak condition number]\label{rem:weak-condition}
For absolutely continuous distributions, the probability that a random perturbation of size exactly $\Delta(x)$ leads to misclassification is zero. For a random direction $v$ uniformly distributed on $\mathbb{S}^{n-1}$, define the directional distance $\Delta(x;v) = \inf\{\eta > 0 \colon \hat{j}(x+\eta v) \neq \hat{j}(x)\}$. If $C_{\hat{j}(x)}$ is star-shaped with respect to~$x$, the $(1-\delta)$-quantile of $\|x\|/\Delta(x;v)$ defines a \emph{weak condition number}
\begin{equation*}
  \mC_\delta(x) := \inf\left\{\eta > 0 \colon \Prob_{v\in \mathbb{S}^{n-1}}\!\left(\frac{\|x\|}{\Delta(x;v)} > \eta\right) \leq \delta\right\},
\end{equation*}
in analogy with the weak condition numbers introduced in~\cite{lotz2020weak} and further studied in~\cite{burgisser2013condition}. The weak condition number captures the typical-case rather than worst-case sensitivity of a classification to perturbations.
\end{remark}


\subsection{Tail bounds on the condition number}
We are interested in the probability that the condition number $\mC(X)$ exceeds a threshold~$t$, for $X$ drawn from a probability distribution on the data space. This is an instance of the probabilistic analysis of condition numbers, studied in detail in~\cite{burgisser2013condition}. In what follows, we distinguish two settings:
\begin{enumerate}
\item \emph{Random data:} $X$ is drawn from a distribution on the data space (average-case analysis);
\item \emph{Random perturbation:} $X = \overline{x} + \xi$ for a fixed data point $\overline{x}$ and a random perturbation $\xi$ (smoothed analysis).
\end{enumerate}
The tube bounds of Section~\ref{sec:tubular} apply to both settings.

\begin{theorem}[Uniform tail bound]\label{thm:condition-uniform}
Let $F\colon \R^n \to \R^m$ be a classifier whose pairwise decision boundaries $V_{ij} = \mathcal{Z}(g_{ij})$ are bounded Pfaffian hypersurfaces with format $(\alpha, \beta, s)$, $s\geq n$, and non-vanishing gradient. Let $X$ be uniformly distributed in $B(p, \sigma)$. Then for $t > 0$,
\begin{equation*}
  \Prob(\mC_p(X) > t) \leq \binom{m}{2}\, C_{\alpha,\beta,s,n}\left(1 + \frac{\alpha + \overline{\beta} + 1}{t}\right)^n,
\end{equation*}
where $C_{\alpha,\beta,s,n}$ and $\overline{\beta}$ are as in Theorem~\ref{thm:prob_bound}.
\end{theorem}

\begin{proof}
  For $X\in B(p,\sigma)$ we have $\|X - p\|\leq \sigma$, so the event $\mC_p(X) > t$ implies $\Delta(X) < \|X-p\|/t \leq \sigma/t$. Since $\Sigma = \bigcup_{i<j}V_{ij}$, we get $\{\Delta(X) < \sigma/t\} \subseteq \bigcup_{i<j}\{\dist(X, V_{ij}) < \sigma/t\}$. By a union bound and Theorem~\ref{thm:prob_bound} applied to each $V_{ij}$ with $\varepsilon = \sigma/t$,
\begin{equation*}
  \Prob(\mC_p(X) > t) \leq \binom{m}{2}\, C_{\alpha,\beta,s,n}\left(1 + (\alpha + \overline{\beta} + 1)\frac{1}{t}\right)^n. \qedhere
\end{equation*}
\end{proof}

\begin{remark}\label{rem:global-condition}
For the global condition number $\mC(X) = \|X\|/\Delta(X)$, the bound $\|X\|\leq \|p\| + \sigma$ gives the variant
\begin{equation*}
  \Prob(\mC(X) > t) \leq \binom{m}{2}\, C_{\alpha,\beta,s,n}\left(1 + (\alpha + \overline{\beta} + 1)\frac{\|p\| + \sigma}{t\sigma}\right)^n.
\end{equation*}
The local form of Theorem~\ref{thm:condition-uniform} is preferred because it is independent of~$\sigma$ and has $O(1/t)$ tail decay for large~$t$.
\end{remark}

We now turn to the smoothed analysis setting, where the data point is a Gaussian perturbation of a fixed input $\overline{x}$. The following result is sharper than the uniform bound, as the Gaussian concentration around $\overline{x}$ yields a tighter estimate via a Fubini decomposition.

\begin{theorem}[Gaussian tail bound]\label{thm:condition-gaussian}
Under the hypotheses of Theorem~\ref{thm:condition-uniform}, let $X$ be normally distributed around $\overline{x}\in \R^n$ with covariance $\sigma^2 \mathrm{Id}$. Then for $t > 0$,
\begin{equation*}
\Prob(\mC_{\overline{x}}(X) > t) \leq \binom{m}{2}\, C_{\alpha,\beta,s,n}\left[\left(1+\frac{\alpha+\overline{\beta}+1}{t}\right)^n - \left(1+\frac{1}{t}\right)^n\right].
\end{equation*}
\end{theorem}

\begin{remark}
For large $t$, the Gaussian bound decays as $\binom{m}{2}\, C_{\alpha,\beta,s,n}\cdot n(\alpha+\overline{\beta})/t$.
\end{remark}

The proof of Theorem~\ref{thm:condition-gaussian} relies on a standard reduction from Gaussian to uniform distributions (cf.~\cite[Section~3.1]{burgisser2013condition}).

\begin{lemma}\label{lem:gauss-to-uniform}
Let $V = \mathcal{Z}(f)$ be a bounded Pfaffian hypersurface with format $(\alpha, \beta, s)$, $s\geq n$, and non-vanishing gradient. Let $X$ be normally distributed around $\overline{x}\in \R^{n}$ with covariance $\sigma^2 \mathrm{Id}$. Then
\begin{equation*}
\Prob(\|X-\overline{x}\|\cdot \dist(X,V)^{-1}\geq t) \leq C_{\alpha,\beta,s,n}\left[\left(1+\frac{\alpha+\overline{\beta}+1}{t}\right)^n - \left(1+\frac{1}{t}\right)^n\right].
\end{equation*}
\end{lemma}

\begin{proof}
Set $A:=\{x\in \R^{n} \mid \|x-\overline{x}\|\cdot \dist(x,V)^{-1}\geq t\}$.
Let $U_{s}$ be uniformly distributed on a closed ball $B(\overline{x},s)$ around $\overline{x}$ of radius $s$, with density
\begin{equation*}
  \frac{1}{\omega_n s^{n}}\cdot \onemtx\{\|x-\overline{x}\|\leq s\},
\end{equation*}
where $\omega_n = \vol B(0,1)$.
For a Gaussian vector $X$ centred at $\overline{x}$ with covariance $\sigma^2\mathrm{Id}$, a standard Fubini argument (see, e.g.,~\cite{burgisser2013condition}) gives
\begin{equation}\label{eq:gauss-to-uniform}
 \Prob(X\in A) = \frac{\omega_n}{(2\pi)^{\frac{n}{2}}}\int_{0}^{\infty} \Prob(U_{\sigma s}\in A)\, s^{n+1} \econst^{-\frac{s^2}{2}} \ \diff{s}.
\end{equation}
Since $U_{\sigma s}\in A$ implies $\dist(U_{\sigma s},V)\leq \sigma s/t$, we can apply Corollary~\ref{cor:tight_prob_bound} with ball radius $\sigma s$ and $\varepsilon = \sigma s/t$. The ratio $\varepsilon/(\sigma s) = 1/t$ is independent of $s$, so
\begin{equation*}
  \Prob(U_{\sigma s}\in A) \leq C_{\alpha,\beta,s,n}\left[\left(1+\frac{\alpha+\overline{\beta}+1}{t}\right)^n - \left(1+\frac{1}{t}\right)^n\right].
\end{equation*}
This is independent of $s$ and factors out of the integral~\eqref{eq:gauss-to-uniform}. The remaining integral evaluates to
\begin{equation*}
  \frac{\omega_n}{(2\pi)^{\frac{n}{2}}}\int_{0}^{\infty} s^{n+1} \econst^{-\frac{s^2}{2}} \ \diff{s} = 1,
\end{equation*}
which follows from $\omega_n = \pi^{n/2}/\Gamma(n/2+1)$ and the substitution $u=s^2/2$ in the Gamma function.
\end{proof}

\begin{proof}[Proof of Theorem~\ref{thm:condition-gaussian}]
The event $\mC_{\overline{x}}(X) > t$ implies $\|X - \overline{x}\|/\dist(X,\Sigma) > t$, which in turn implies $\|X-\overline{x}\|/\dist(X,V_{ij}) > t$ for some pair $i < j$. By a union bound and Lemma~\ref{lem:gauss-to-uniform} applied to each~$V_{ij}$,
\begin{equation*}
  \Prob(\mC_{\overline{x}}(X) > t) \leq \binom{m}{2}\, C_{\alpha,\beta,s,n}\left[\left(1+\frac{\alpha+\overline{\beta}+1}{t}\right)^n - \left(1+\frac{1}{t}\right)^n\right]. \qedhere
\end{equation*}
\end{proof}


\subsection{Neural network classifiers}
We now specialise the tail bounds to neural network classifiers with Pfaffian activation functions.

\begin{corollary}[Sigmoid/tanh network]\label{cor:nn-condition}
Let $F\colon \R^n\to \R^m$ be a fully connected neural network classifier with $\ell$ hidden layers, $h$ total hidden units, and an activation function with Pfaffian format $(\alpha, \beta, s)$ that is autonomous with $s\geq 1$. Assume the pairwise decision boundaries $V_{ij}$ are bounded with non-vanishing gradient.

If $X$ is uniformly distributed in $B(p,\sigma)$, then
\begin{equation*}
  \Prob(\mC_p(X) > t) \leq \binom{m}{2}\, C_{\alpha_F,\beta_F,s_F,n}\left(1 + \frac{\alpha_F + \overline{\beta}_F + 1}{t}\right)^n,
\end{equation*}
where $\alpha_F = \ell(\alpha+\beta-1)-\beta+1$, $\beta_F = \beta$, $s_F = sh$, and $\overline{\beta}_F = \max\{\beta_F, 2\}$.
\end{corollary}

\begin{proof}
By Proposition~\ref{prop:nn-format}, the network output has Pfaffian format $(\alpha_F, \beta_F, s_F)$. Since $g_{ij} = F_j - F_i$ shares the same Pfaffian chain (Remark~\ref{re:1}), each $V_{ij}$ is a Pfaffian hypersurface with the same format. The result follows from Theorem~\ref{thm:condition-uniform}.
\end{proof}

\begin{example}[Sigmoid activation]\label{ex:sigmoid-condition}
For the logistic sigmoid $\sigma(x) = (1+\econst^{-x})^{-1}$ with format $(2,1,1)$, we get $\alpha_F = 2\ell$, $\beta_F = 1$, $s_F = h$, and $\overline{\beta}_F = 2$, so
\begin{equation*}
  \Prob(\mC_p(X) > t) \leq \binom{m}{2}\cdot 6\cdot 2^{\frac{h(h-1)}{2}}\big(n(4\ell+1)+2\big)^h\left(1+\frac{2\ell+3}{t}\right)^n,
\end{equation*}
provided $h\geq n$. The Gaussian variant (Theorem~\ref{thm:condition-gaussian}) gives, for $X\sim \mathcal{N}(\overline{x}, \sigma^2\mathrm{Id})$,
\begin{equation*}
  \Prob(\mC_{\overline{x}}(X) > t) \leq \binom{m}{2}\cdot 6\cdot 2^{\frac{h(h-1)}{2}}\big(n(4\ell+1)+2\big)^h\left[\left(1+\frac{2\ell+3}{t}\right)^n - \left(1+\frac{1}{t}\right)^n\right].
\end{equation*}
For large $t$, the Gaussian bound decays as $\binom{m}{2}\cdot 6\cdot 2^{h(h-1)/2}(n(4\ell+1)+2)^h \cdot n(2\ell+2)/t$.
\end{example}

\begin{remark}\label{rem:bound-interpretation}
The dependence on the network architecture enters through two terms: the Khovanskii prefactor $2^{h(h-1)/2}(n(4\ell+1)+2)^h$, which grows rapidly with the number of hidden units~$h$, and the polynomial tail $(1+(2\ell+3)/t)^n$, which depends on the depth~$\ell$ and the input dimension~$n$. For practical networks, the Khovanskii prefactor renders the bound quantitatively vacuous; however, the bound is qualitatively meaningful in several ways:
\begin{enumerate}
\item The $O(1/t)$ tail decay is sharp in the sense that it matches the generic behaviour of condition number tails for smooth hypersurfaces;
\item The bound captures an explicit depth-width tradeoff: deeper networks increase the polynomial base while wider networks increase the exponential prefactor;
\item For fixed architecture, the bound gives the correct scaling in $\varepsilon$ and~$n$.
\end{enumerate}
\end{remark}
